\section{Schriftarten: \datei{OTHR\_fonts.sty}}

Die Hausschrift des Corporate Designs der OTH Regensburg ist die Lucida Sans.
Um diesem Schriftbild möglichst nah zu kommen, unterstützt das Zeichensatzpaket verschiedene Optionen.

Default ist die Verwendung der DejaVU Sans, welche dem Original sehr ähnlich ist.
Diese Schrift wird bei vielen Distributionen im Standardumfang mitgeliefert, bietet aber in der Type~1 Variante keinen Zeichensatz für den Formelsatz, der daher mit pdflatex traditionell gesetzt wird.

Bei Verwendung von LuaTeX wird versucht, als Standard eine der Originalschriftarten einzusetzen, sofern diese Auf dem Rechner installiert sind.
Alternativ wird auch hier die OpenType"=Variante von DejaVU Sans inklusive des Zeichensatzes für den Formelsatz verwendet.

Falls eine Lizenz für die Lucida"=Zeichensätze vorhanden ist, kann auch bei pdflatex über die Option \texttt{Lucida}\docoption{Lucida\\lucida} auf die Lucida"=Schriftfamilie umgeschaltet werden, die auch den Formelsatz unterstützt und insgesamt ein sehr harmonisches Schriftbild liefert\footnote{Lizenzen für die Lucida Fonts im OpenType"= und T1"= Format sind bei der \TeX Users Group erhältlich: \href{https://www.tug.org/store/lucida/index.html}{https://www.tug.org/store/lucida/index.html}.}.

Bei Verwendung von LuaTeX und einer vorhandenen Installation der Lucida Schriftart Grande auf dem System (typisch für Apple"=Systeme), kann über die Option \texttt{LucidaGrande}\docoption{LucidaGrande} auch diese Schriftart verwendet werden.

Da sich in manchen Fällen die serifenlose Schriftart nur bedingt eignet (z.B. sehr lange geschlossene Texte mit relativ langen Zeilen bzw. hoher Zeichenanzahl) kann über die Option \texttt{KeepRoman}\docoption{KeepRoman} die Variante der jeweils gewählten Schrift mit Serifen als Standardschrift beibehalten werden.
